
\documentclass{article}

\usepackage{fancyhdr, graphicx, pslatex, array}
\usepackage[table]{xcolor}

\raggedbottom
\widowpenalty=1000
\clubpenalty=1000

\newcommand{\stdfooter}{%
  \fancyhf{}
  \fancyhfoffset{0.2in}
  \fancyfoot[L]{
     \parbox{\logowidth}{
       \includegraphics[width=\logowidth]{labb_logo}\\
       \vskip 0.1in
       \includegraphics[width=\logowidth]{drcs_logo-200}
     }
     \setlength\acklength{\textwidth}
     \addtolength\acklength{-\logowidth}
     \addtolength\acklength{-0.5in}
     \hskip 0.5in
     \parbox{\acklength}{
       This report was generated at www.myairsample.org, a site developed and
       maintained by the Louisiana Bucket Brigade and Digitial Resources for
       Community and Science.  For questions and comments about air data,
       Louisiana Bucket Brigade can contacted at www.labucketbrigade.org,
       4226 Canal St, New Orleans, LA 70119, phone: 504-484-3433, fax:
       504-324-0332.  For comments, feedback, or errors in the web site,
       please contact the LA Bucket Brigade or email Digital Resources for
       Community and Science at drcsdirector@gmail.com
     }
  }
}

\fancypagestyle{plain}{\stdfooter}
\pagestyle{empty}
% \pagestyle{fancy}
\stdfooter

\setlength{\parindent}{0pt}%
\setlength{\parskip}{\baselineskip}%
\setlength{\oddsidemargin}{0in}
\setlength{\textwidth}{6.5in}
\newlength\acklength
\newlength\logowidth
\setlength\logowidth{1.8in}
\addtolength\voffset{-0.5in}

\setlength{\arrayrulewidth}{1.5pt}

\newcommand{\cubed}{$^3$}
\newcommand{\micro}{\ensuremath{\mu}}
\newcommand{\fc}{\cellcolor{salmon}}
\newcommand{\nfc}{\cellcolor{white}}
\definecolor{salmon}{rgb}{1.0,0.8,0.7}
\newcommand{\highlightbox}[1]{\colorbox{salmon}{\parbox{\linewidth}{#1}}}

\newcommand{\inunits}        {\subst{inunits}} 
\newcommand{\outunits}       {\subst{outunits}} 
\newcommand{\longinunits}    {\subst{longinunits}} 
\newcommand{\longoutunits}   {\subst{longoutunits}} 
\newcommand{\samplename}     {\subst{samplename}} 
\newcommand{\sampledate}     {\subst{sampledate}} 
\newcommand{\samplelocation} {\subst{samplelocation}} 
\newcommand{\user}           {\subst{user}}


\begin{document}

\title{Levels of Concern Report}
\author{\user}
\date{}
\maketitle

\section*{Sample information}
\begin{itemize}
\ifx\samplename\undefined\relax\else\item Sample identifying name:   \samplename\fi
\ifx\sampledate\undefined\relax\else\item Date sample was taken:     \sampledate\fi
\ifx\samplelocation\undefined\relax\else\item Location sample was taken: \samplelocation\fi
\item Report date:               \today
\item Report input units:        \longinunits\ (\inunits)$^*$
\item Report output units:       \longoutunits\ (\outunits)$^*$
\item Report made on web at:     http://myairsample.org
\end{itemize}

$^*$For a description of what units mean, see ``Units Information'' section later in this report

\newcommand{\unitssection}{
\section*{Units information}
Parts per billion (ppb) describes how many weighed parts of a chemical there
are for 1 billion parts of air. For example, a recipe says to add just a drop of
vanilla for every 100 pounds of flour. The drop of vanilla weighs hardly
anything, but it has a big effect on the cookies' flavor. Similarly, if we measure
benzene in the air, we might find 3 ``drops'' of benzene for 1,000,000,000
(one billion) ``drops'' of air. It seems like a small amount, but it is significant.

Parts per billion by volume, or ppbv, means the concentration has been figured
out in terms of how much space the molecules take up. For example, if we
make a mixture of 3 cups of vanilla and 1 billion cups of flour, then our
concentration is 3 parts volume (cups of vanilla) per billion parts volume (cups
of flour), or 3 ppbv sugar in flour. When 3 volumes of benzene are in a billion
volumes of air, the concentration is 3 ppbv benzene in air.

Micrograms per meters cubed (\micro g/m\cubed) describes how much of a chemical's
weight is in a volume of air that takes up one cubic meter. Imagine an empty
box that is three feet long on both sides, and three feet tall. One meter is about
three feet long. So the box's volume is 1 cubic meter, or 1 m\cubed. A microgram
(\micro g) is a very small weight, like that of a grain of sand. You put 3 grains of sand
into the box. The concentration of sand inside the box is the weight of the sand
(3 \micro g) divided by the
volume of the box (1 m\cubed), or 3 \micro g /m\cubed. Like grains of sand, chemicals can also be
reported by weight and volume. For example, a monitor might read 5 \micro g /m\cubed\ %
benzene, or 5 \micro g of benzene in 1 m\cubed\ of air.
}

\subst{resultssection}

\unitssection

\section*{Sample screening levels}

\highlightbox{Some government agencies have developed standards and screening levels for
toxic chemicals in the air based on health information about the chemicals.
There is no information available for some toxic chemicals. The agencies are
listed below, with a brief description of the methods used in establishing their
levels. States may not be required to adhere to national standards.}

\begin{itemize}

\item Texas Effects Screening Levels


http://www.tceq.texas.gov/toxicology/esl


These levels are based on existing studies of chemical health
effects. Below these levels, no adverse health effects are
thought likely to occur. They reflect the
experimentally-determined levels at which the chemicals caused
adverse effects in study populations (NOAELs or LOAELs) of people
or animals, combined with safety factors to account for the
differences among human populations and between humans and
animals.


"Short-term" refers to an exposure duration of one hour. "Long-term"
refers in most cases to a duration of one year; for benzene and
ethylene dichloride, it indicates a 24-hour period.


These levels reflect both cancer and non-cancer effects. They are not
legally enforceable.
\item Louisiana Ambient Air Quality Standards


http://www.deq.louisiana.gov/portal/tabid/1674/Default.aspx


These levels are legally enforceable standards in Louisiana, developed
through Louisiana's regulatory process. They are found in Table 51.2 of
Title 33, Part III.
They are based on health effects information about the chemicals: the
eight-hour standard modifies occupational exposure levels to be
appropriate for residential exposures; the annual standard is based on
EPA procedures for calculating cancer risks.
\item ATSDR (American Toxic Substance and Disease Registry) Minimal Risk Levels


http://www.atsdr.cdc.gov/mrls/index.asp


These levels are based on existing studies of chemical health
effects. Exposure at these levels is said to pose "minimal risks"
of adverse health effects. The levels were derived from the
experimentally-determined levels at which the chemicals caused
adverse effects in study populations (NOAELs or LOAELs) of people
or animals, combined with uncertainty factors to account for the
differences among human populations and between humans and
animals.


"Acute" exposure is defined as up to two weeks, "Intermediate" as
between two weeks and one year, and "Chronic" as longer than a year.


These levels reflect only non-cancer health effects. They are not
legally enforceable.
\item North Carolina Ambient Air Standards


http://daq.state.nc.us/Rules/rules/secD1100.pdf


These levels are legally enforceable standards in North Carolina, developed through North Carolina's regulatory process.


They are based on health effects information about the chemicals.
\item EPA Region 6 Screening Levels


http://www.epa.gov/earth1r6/6pd/rcra/pd-n/screen.htm


These levels are based on existing studies of chemical health effects.
They levels are calculated for residential (as opposed to workplace)
exposures. They reflect the risks of exposure to a certain level of the
chemical. The levels listed as screening levels correspond to pre-
determined levels of risk from exposure: either 1 in a million cancer risk
or a "hazard quotient" of 1 for non-cancer effects, whichever
corresponds to a lower concentration. These screening levels are not
legally enforceable.

\end{itemize}


\end{document}

